%; whizzy paragraph -pdf xpdf -latex ./whizzypdfptex.sh
%; whizzy-paragraph "^\\\\begin{frame}\\|\\\\emtext"
% latex beamer presentation.
% platex, latex-beamer でコンパイルすることを想定。 

%     Tokyo Debian Meeting resources
%     Copyright (C) 2012 Junichi Uekawa
%     Copyright (C) 2011 Nobuhiro Iwamatsu

%     This program is free software; you can redistribute it and/or modify
%     it under the terms of the GNU General Public License as published by
%     the Free Software Foundation; either version 2 of the License, or
%     (at your option) any later version.

%     This program is distributed in the hope that it will be useful,
%     but WITHOUT ANY WARRANTY; without even the implied warreanty of
%     MERCHANTABILITY or FITNESS FOR A PARTICULAR PURPOSE.  See the
%     GNU General Public License for more details.

%     You should have received a copy of the GNU General Public License
%     along with this program; if not, write to the Free Software
%     Foundation, Inc., 51 Franklin St, Fifth Floor, Boston, MA  02110-1301 USA

\documentclass[cjk,dvipdfmx,12pt]{beamer}
\usetheme{Tokyo}
\usepackage{monthlypresentation}

%  preview (shell-command (concat "evince " (replace-regexp-in-string "tex$" "pdf"(buffer-file-name)) "&")) 
%  presentation (shell-command (concat "xpdf -fullscreen " (replace-regexp-in-string "tex$" "pdf"(buffer-file-name)) "&"))
%  presentation (shell-command (concat "evince " (replace-regexp-in-string "tex$" "pdf"(buffer-file-name)) "&"))

%http://www.naney.org/diki/dk/hyperref.html
%日本語EUC系環境の時
\AtBeginDvi{\special{pdf:tounicode EUC-UCS2}}
%シフトJIS系環境の時
%\AtBeginDvi{\special{pdf:tounicode 90ms-RKSJ-UCS2}}

\title{東京エリアDebian勉強会}
\subtitle{第88回 2012年5月度}
\author{上川純一\\dancer@debian.org}
\date{2012年5月19日}
\logo{\includegraphics[width=8cm]{image200607/openlogo-light.eps}}

\begin{document}

\frame{\titlepage{}}

\begin{frame}{設営準備にご協力ください。}
会場設営よろしくおねがいします。
\end{frame}

\begin{frame}{Agenda}
\begin{minipage}[t]{0.45\hsize}
  \begin{itemize}
  \item 注意事項
	\begin{itemize}
	 \item 飲食禁止
	 \item 宗教禁止
	 \item 営利活動禁止
	\end{itemize}
   \item 最近あったDebian関連のイベント報告
	\begin{itemize}
        \item 第87回 東京エリアDebian勉強会
	\end{itemize}
 \end{itemize}
\end{minipage} 
\begin{minipage}[t]{0.45\hsize}
 \begin{itemize}
  \item Debian Trivia Quiz
  \item Python初心者が「Pythonプロフェッショナルプログラミング」を読んでみた
  \item coffeescriptをDebianで使う
  \item 月刊Debhelper 
 \end{itemize}
\end{minipage}
\end{frame}

\section{DWN quiz}
\emtext{DWN quiz}
\begin{frame}{Debian 常識クイズ}

Debian の常識、もちろん知ってますよね?
知らないなんて恥ずかしくて、知らないとは言えないあんなことやこんなこと、
みんなで確認してみましょう。

今回の出題範囲は\url{debian-devel-announce@lists.debian.org},
\url{debian-devel@lists.debian.org} に投稿された
内容とDebian Project Newsなどからです。

\end{frame}

\subsection{問題}
%; whizzy-master ../debianmeetingresume201205.tex
% 以上の設定をしているため、このファイルで M-x whizzytex すると、whizzytexが利用できます。
%

\santaku
{Debian installer 7.0 alpha 1 のリリース日は}
{5/13}
{6/13}
{4/13}
{A}
{Wheezyのインストーラーのアルファリリースが出たので皆さん試してください。}

\santaku
{Cyril Brulebが6月にWheezyをフリーズすると発表したが、Transitionの締め切りはいつだといっているか}
{5月13日}
{6月10日}
{5月20日}
{C}
{Transitionするなら5月20日までにバグをファイルしておけとのこと。Transitionというのはざっくりというと多数のパッケージが相互に依存しているような変更。例えば、ライブラリのABIが変わるだとか。}


\section{事前課題}
\emtext{事前課題}
{\footnotesize
 %; whizzy-master ../debianmeetingresume201205.tex
% $B0J>e$N@_Dj$r$7$F$$$k$?$a!"$3$N%U%!%$%k$G(B M-x whizzytex $B$9$k$H!"(Bwhizzytex$B$,MxMQ$G$-$^$9!#(B
%

\begin{prework}{ BeatenAvenue }

1.$B$*$9$9$aK\!'(B

$B!V(Bwindows$B%W%m%U%'%C%7%g%J%k%2!<%`%W%m%0%i%_%s%0(B($BA4#24,(B?)$B!W(B
DirectX$B4XO"$NFbMF$b$"$C$?5$$,$7$^$9$,(Bwin32API$B$K$D$$$F$NFbMF$,B?$+$C$?5$$,$7$^$9!#@N$NK\$G$9$,%?%9%/=hM}$N9M$(J}$J$I8=:_$G$bDLMQ$9$kItJ,$OB?$$$H;W$$$^$9!#$3$NK\$rGc$C$F$+$i(BC++$B$NJY6/$r;O$a$?$3$H$b$"$C$F!&!&!&8D?ME*$K;W$$=P$,$?$/$5$s$"$j$^$9!#(B

$B!V(BDirectX$B5U0z$-BgA4(B500$B$N6K0U!W(B
$BF~LgE*$JM%$7$$2r@b$+$i0lJbF'$_9~$s$@%W%i%9%"%k%U%!$^$GB7$C$F$$$^$9!#;DG0$J$,$i@dHG$G?^=q4[$+$i<Z$j$FFI$_$^$7$?!#(BDirectX9$B$N2r@b=q$G$O0lHV$h$$$b$N$+$H;W$$$^$9!#(B

$B!V(BGameProgrammingGems($B%7%j!<%:(B)$B!W(B
$B3$30$N%2!<%`%W%m%0%i%^$NJ}!9$,=q$$$?5-;v$r$^$H$a$?K\!#(B3$B4,$@$C$?$+$H;W$$(B
 $B$^$9$,(BNaughty Dog$B$NJ}$,=q$$$?(B"Jak and Daxter: The Precursor legacy"$B$G$N%^%C%W0\F0=hM}$K$D$$$F$NFbMF$,9%$-$G$9!#$*CMCJ0J>e!#(B

2.$B=i?4<T$K$*$9$9$a$9$k%9%/%j%W%H8@8l!'(B
Debian$B$H$OA4$/4X78$J$$$G$9$,(BDOS$B%P%C%A%U%!%$%k$H(BExcelVBA$B$N$A$g$C$H$7$?;H$$J}$O3P$($J$$$H;vL3;E;v$,?J$_$^$;$s!#2r@b%5%$%H$bB?$$$N$GIU$-$C$-$j$G65$($kI,MW$b$"$^$j$J$$$G$9!#(B
Linux$B$@$H(Bbash$B$N%9%/%j%W%H$J$s$G$7$g$&$+!#;d$,=i?4<T$J$N$G$=$l$7$+?($C$F$$$^$;$s!&!&!&!#(B
\end{prework}

\begin{prework}{ amotoki }

1. $B:rF|M'?M$+$i4+$a$i$l$F5$$K$J$C$F$$$k$N$,!V>pG.%W%i%0%i%^!<!W$G$9!#<+J,$N?M@8$r<+J,$G@Z$j3+$$$F$$$/$?$a$KI,MW$J$3$H$,J,$+$j$d$9$/@0M}$5$l$F$$$k$1$I!"<+J,$rA0$K?J$a$F$/$l$k>pG.$r46$8$?$H$N$3$H!#$5$C$=$/CmJ8$7$?!#(B

2. $B:#$*$9$9$a$9$k$H$7$?$i(BPython$B$r$*A&$a$7$^$9!#%*%V%8%'%/%H;X8~$b=q$-$d$9$$$7!"%i%$%V%i%j$b=<<B$7$F$$$F!"%^%K%e%"%k$b<BNc$,$=$m$C$F$$$k$N$G!"%W%m%0%i%_%s%0$r3X$s$G$$$/>e$G$h$$$H;W$$$^$9!#Bg$-$a$N(BOSS$B%W%m%8%'%/%H$G$b;H$o$l$F$$$k$N$G!"CN$C$F$*$$$FB;$9$k$3$H$b$"$j$^$;$s!#:G=i$N0lJb$H$7$F$O!V(BPython$B%A%e!<%H%j%"%k!W$,$h$$$H;W$$$^$9!#:#$G$b$H$-$I$-8+$k$3$H$,$"$j$^$9!#(B
\end{prework}

\begin{prework}{ $B5HLn(B(yy\_y\_ja\_jp) }

1. Git$B$K$h$k%P!<%8%g%s4IM}(B $B<B:]$N%W%m%8%'%/%H$G$N(BGit$B$NMxMQK!$,=q$+$l$F$$$k$h$&$G$9!%(B

2. $B%7%'%k%9%/%j%W%H(B
$B5$7Z$K;H$($k$+$i$G$9!%(B
bash(1), dash(1)
\end{prework}

\begin{prework}{ dictoss($B?yK\!!E5=<(B) }

1.$B!V%$%s%F%k(B $B%9%l%C%G%#%s%0!&%S%k%G%#%s%0!&%V%m%C%/(B -- $B%^%k%A%3%";~Be$N(BC++$BJBNs%W%m%0%i%_%s%0!W%*%i%$%j!<!&%8%c%Q%s!"(BJames Reinders$BCx(B
Intel$B$,3+H/$78=:_$O(BGPLv2$B$G8x3+$7$F$$$k(BC++$B$NJBNs7W;;MQ%i%$%V%i%j!V(BThreading Building Blocks$B!W$N2r@b$r9T$C$F$$$kK\!#%^%k%A%3%";~Be$NCf$GJ#?t$N(BCPU$B%3%"$r8zN(E*$K;HMQ$7$F7W;;@-G=$r>e$2$k$?$a$NCN<1$,5M$^$C$F$$$k!#$"$/$^$GC10l%N!<%I$G7W;;@-G=$r>e$2$k$?$a$N<jK!$G$"$j!"J#?t%N!<%I$G7W;;@-G=$r8~>e$5$;$k%/%i%9%?%j%s%05;=Q$NOC$G$O$J$$$N$GCm0U!#(B

2.$B$*4+$a$O(Bpython$B!#M}M3$O%W%m%0%i%_%s%0=i?4<T$H$$$&$3$H$G3+H/<T$K$h$C$F=q$-J}$K:90[$,=P$K$/$$J,(Bweb$B$G>R2p$5$l$F$$$k%3!<%I$KJJ$,$J$/$H$C$D$-$d$9$$$?$a!#(B
$B$*$9$9$a%5%$%H$O!"$&!<$s!"$J$s$G$7$g!)<+J,$OJL$N8@8l$,=q$1$k$h$&$K$J$C$F$+$i(Bpython$B$r;O$a$?$N$G(B"\url{http://www.python.jp/doc/release/}"$B$r3NG'$7$^$9!#(B
\end{prework}

\begin{prework}{ kamonshohei }

1. $B%7%'%k%9%/%j%W%H%7%'%k%9%/%j%W%H4pK\%j%U%!%l%s%9(B

2. $B%7%'%k%9%/%j%W%H$G$7$g$&$+!#%j%J%C%/%9$N%3%^%s%I$@$1$G!"$A$c$A$c$C$H<BAu$G$-$k<j7Z$5$,$$$$$G$9!#:G=i$N0lJb$G0FFb$9$k=q@R$O(B 1$B$G$"$2$?%7%'%k%9%/%j%W%H4pK\%l%U%!%l%s%9$G$9!#(B
\end{prework}

\begin{prework}{ emasaka }

1. $B%1%s!&%9%_%9!VC/$b65$($F$/$l$J$$@;=q$NFI$_J}!W!#@;=q$K=q$+$l$F$$$k$=$N$^$^$NJ8LL$r??LLL\$KFI$s$G%f!<%b%i%9$K>R2p!J!)!K$7$F$$$kK\!J(BDebian$B4X78$J$$!K(B

2. Bash$B$H=q$3$&$+$H;W$C$?$1$I(BRuby$B!#M}M3$O!V%*%V%8%'%/%H;X8~$r;H$C$F$b;H$o$J$/$F$b$$$$!W$G$O$J$/$F!V%*%V%8%'%/%H;X8~$r6/@)$5$l$k!W$+$i!#=q@R$O!V$?$N$7$$(BRuby$B!W(B

\end{prework}

\begin{prework}{ $BK\>1(B }

1. Debian $BJY6/2q;22C<T$K>R2p$7$?$$=q@R$r#1:}0J>e5s$2$F!"FbMF$r4JC1$K>R2p$7$F$/$@$5$$!#(B
$BDjHV$G$9$,!X%O%C%+!<$H2h2H!Y$H$+$I$&$G$7$g$&!#%*%?%/$N1eG0$,9~$a$i$l$F$$$^$9!#?F$H$7$F;R6!$N>-Mh$KIT0B$r46$8$^$9!#(B

2. $B$"$J$?$,2?$+%9%/%j%W%H8@8l$r%W%m%0%i%_%s%0=i?4<T$K$*4+$a$9$k$H$7$F!V$=$N8@8l$rA*$s$@M}M3!W$H!V:G=i$N0lJb$H$7$F0FFb$9$k=q@R!?%5%$%H!W$r65$($F$/$@$5$$!#(B
PHP$B$,$*$9$9$a$G$9!#$[$+$KHf$Y$F;E;v$,B?$=$&$H$$$&M}M3$G$9!#B?$$$+$I$&$+$O<B:]$N$H$3$m$o$+$j$^$;$s$,!"(B

\url{http://www.google.co.jp/trends/?q=PHP,+Perl,+Python,+Ruby,+Javascript,+Haskell&ctab=0&geo=jp&geor=all&date=ytd&sort=0}

$B$3$38+$k$HB?$=$&$G$9!#0FFb$9$k=q@R$H$7$FDjHV$O%^%s%b%9K\$@$H;W$$$^$9$,!"FI$s$@$3$H$O$"$j$^$;$s!#$=$7$F8E$$>pJs$+$b$7$l$^$;$s!#(B
$B0JA0!"$H$"$k(BPHP$BJ}LL$NJ}$H$*OC$7$9$k5!2q$,$"$j!"%*%i%$%j!<$NK\$O$A$g$C$H!DE*$J$3$H$rOC$7$?$i!"!V$"$"$$$&K\$b$$$$$+$J$H;W$C$F$^$9!W$H$$$o$l$^$7$?!#K]Lu<T$G$7$?!#(B

\end{prework}

\begin{prework}{ henrich }

1. $B!VF~Lg(BDebian$B%Q%C%1!<%8!W!#=qL>$+$iFbMF$,J,$+$k$+$H$O;W$$$^$9$,!"(BDebian$B%Q%C%1!<%8$N:n$jJ}$N=q@R$G$9!#B34)$,4|BT$5$l$^$9(B

2. Python $B$rA*$S$^$7$?!#(BPerl$B$O?M$K$h$C$F$H$F$b=q$-J}$,JQ$o$k$H$3$m$,$"$^$j4r$7$/$J$/!"(BRuby$B$O%P!<%8%g%s4V$N0\9T$,<c43MpK=$K46$8$i$l$?$N$G!#2?EY$b%W%m%0%i%_%s%0$K:C@^$7$F$$$k;d$G$9$,!":#2s(BPython$B$r3X$V$N$KA*$s$@!V=i$a$F$N%3%s%T%e!<%?%5%$%(%s%9!W$,$H$F$bNI$$=q@R$G$7$?!#(B
\end{prework}

\begin{prework}{ $BLnEg!!5.1Q(B }

1.$B!V%$%N%Y!<%7%g%s$N%8%l%s%^!W!J(BISBN10:4798100234)$B$H!"!V$N$&$@$^!W(B(ISBN10:4344015959)$B!#!V%$%N%Y!<%7%g%s$N%8%l%s%^!W$O!"5;=Q3W?7$,4{B8$N$b$N$r$V$A2u$7$F$$$/2aDx$K$*$$$F!"4{B85;=Q$K$*$$$FM%=($JAH?%$G$"$l$P$"$k$[$I5;=Q3W?7$K$D$$$F$$$1$J$/$J$C$F$7$^$&8=>]$rM}5M$a$G@bL@$7$?K\!#!V$N$&$@$^!W$O?M$N9TF0$K$*$$$F$O!"<B$O=,47$,@h$G$d$k5$$O8e$+$i$D$$$F$/$k$b$N$G$"$k$H$$$&;v$r@bL@$7$?K\!#(B

2. $B%W%m%0%i%_%s%0=i?4<T$K$O:#;~$N>u67$+$i(Bjavascript/HTML5$B$r4+$a$?$$$N$G$9$,!"4N?4$N<+J,$,L$I>2A!#:G=i$N0lJb$O%2!<%`M7$S$?$5$K%W%m%0%i%`3P$($?7P83$r85$K!"(B\url{http://wise9.jp/}$B!!$H!"!!(B\url{http://enchantjs.com/}$B!!$,$*$9$9$a$J$N$+$J!)(B
\end{prework}

\begin{prework}{ yamamoto }

1. Debian $BJY6/2q;22C<T$K>R2p$7$?$$=q@R$r#1:}0J>e5s$2$F!"FbMF$r4JC1$K>R2p$7$F$/$@$5$$!JFC$K5;=Q=q$K$O8B$j$^$;$s!K!#(B

$BM-L>$@$H;W$&$N$G!">R2p$9$k$[$I$N$3$H$O$J$$$+$b$7$l$^$;$s$,!"!VF~Lg(BUNIX$B%7%'%k%W%m%0%i%_%s%0(B-$B%7%'%k$N4pAC$+$i3X$V(BUNIX$B$N@$3&!!(BBruce Blinn $BCx!&;32<E/E5(B $BLu!W$r0&FI$7$F$$$^$9!#$^$"!"%7%'%k%9%/%j%W%H%W%m%0%i%_%s%0$r%7%'%k%W%m%0%i%_%s%0$H8@$C$F$$$k%?%$%H%k(B ($BL^O@%7%'%k<+BN$N2r@b$b$"$j$^$9$1$I$M(B) $B$O%"%l$G$9$,!"(BB$B%7%'%k$N%9%/%j%W%H$r=q$/$H$-$K$ONI$/3+$$$F$$$^$9!#(B

2. $B$"$J$?$,2?$+%9%/%j%W%H8@8l$r%W%m%0%i%_%s%0=i?4<T$K$*4+$a$9$k$H$7$F!V$=$N8@8l$rA*$s$@M}M3!W$H!V:G=i$N0lJb$H$7$F0FFb$9$k=q@R!?%5%$%H!W$r65$($F$/$@$5$$!#(B

\#!/bin/sh $BK|:P!*(B
$B$I$3$G$bBg35F0$/$+$i$*$9$9$a!#(B
B$B%7%'%k$N:nK!$K=>$C$F$$$l$P!":#$N(Bdash$B$J$iBgBNF0$/!)$_$?$$(B ($BL$3NG'(B)$B!#(B
$B$@$a$J$i!!(B\#!/bin/bash$B$G!#(B
$B$H$+$$$$$J$,$i$b!"$*$$$i$N%7%'%k%9%/%j%W%H$K$O30It%3%^%s%I%P%j%P%jF~$l$F$^$9$1$I!#(B($B$&$R(B

$B>e5-$N=q@R$G$b$$$$$G$9$,!"$0$0$k$5$s$K$*;G$$$7$?$i!";29M$K$J$k%5%$%H$O;3$[$I=P$F$/$k$G$7$g$&!#(B
\end{prework}

}

\section{python}
\emtext{python}

\section{coffeescript}
\emtext{coffeescript}

 \begin{frame}{Coffeescript}
  \begin{itemize}
   \item javascript を書きやすくするための言語
   \item npmjs.org調べで2番目にポピュラーなモジュール
  \end{itemize} 
\end{frame}

 \begin{frame}{Smooth Coffeescript}
  \url{http://autotelicum.github.com/Smooth-CoffeeScript/}
 \end{frame}

 \begin{frame}[containsverbatim]{install}
  \begin{commandline}
   $ sudo apt-get install coffeescript libjs-coffeescript
  \end{commandline}
 \end{frame}


 \begin{frame}[containsverbatim]{browser}
 \begin{commandline}
<!DOCTYPE HTML PUBLIC "-//W3C//DTD HTML 4.01 Transitional//EN">
<html>
  <head>
    <title>Hello world</title>
    <script type='application/javascript' src='coffee-script.min.js'></script>

    <script type='text/coffeescript' src='./hello.coffee'></script>

    <script type='text/coffeescript' >
console.log 'This is a pen'
    </script>
  </head>
  <body>
    <h1>Hello world</h1>
  </body>
</html>
 \end{commandline}
 \end{frame}

 \begin{frame}[containsverbatim]{client}
  \begin{commandline}
$ cat shebang.coffee
#!/usr/bin/env coffee
util = require 'util'
util.log 'hello world'
$ ./shebang.coffee
8 May 17:08:25 - hello world
  \end{commandline}
 \end{frame}

 \begin{frame}[containsverbatim]{client}
\begin{commandline}
cli = require 'cli'

cli.parse {
    port: ['p', 'port number to listen to', 'number', 8088]
}

http = require 'http'

cli.main (args, options) ->
  http.createServer((req, res) ->
    res.writeHead 200, {'Content-Type': 'text/plain'}
    res.end 'Hello world\n').listen options.port
  console.log 'Server running at http://localhost:' + options.port + '/'
\end{commandline}
 \end{frame}


\section{月刊Debhelper}
\emtext{月刊Debhelper}

\section{今後のイベント}
\emtext{今後のイベント}
\begin{frame}{今後のイベント}
\begin{itemize}
\item 大統一Debian勉強会\\
日時：6月23日(土曜日)\\
場所：京都大学 理学研究科3号館 数学教室 108, 109, 110 号室
\end{itemize}
\end{frame}

\section{今日の宴会場所}
\emtext{今日の宴会場所}
\begin{frame}{今日の宴会場所}
はなの舞 荻窪西口店
\url{http://www.hotpepper.jp/strJ000015000?vos=evpahpp1016x0030058}
\end{frame}

\end{document}

;;; Local Variables: ***
;;; outline-regexp: "\\([ 	]*\\\\\\(documentstyle\\|documentclass\\|emtext\\|section\\|begin{frame}\\)\\*?[ 	]*[[{]\\|[]+\\)" ***
;;; End: ***
